%=====================================================================
\chapter{Three Levels, and Why They Are Separate}\label{ch:architecture}
%=====================================================================
\locator{1}{Part I --- What a Database Is}

\begin{outcomes}
By the end of this chapter you will be able to:
\begin{tight}
  \item \textbf{Name} the three levels of the ANSI/SPARC architecture and
        \textbf{state} what each one describes.
  \item \textbf{Distinguish} a schema from an instance, and \textbf{explain}
        why the difference matters when something changes.
  \item \textbf{Define} logical and physical data independence, and
        \textbf{give} an example of a change that each one protects you from.
  \item \textbf{Identify} which level a given change belongs to.
  \item \textbf{Explain} why a system with no separation of levels is the
        file-based system of \chref{whynotfiles} again.
\end{tight}
\end{outcomes}

\begin{prereq}
\chref{whynotfiles}: the six problems, and what a DBMS is. Nothing else. This
chapter is the one piece of pure architecture in the course and it is here
because every later chapter quietly assumes it.
\end{prereq}

\begin{keyterms}
ANSI/SPARC architecture $\bullet$ external level $\bullet$ external schema
$\bullet$ user view $\bullet$ conceptual level $\bullet$ conceptual schema
$\bullet$ internal level $\bullet$ internal schema $\bullet$ physical level
$\bullet$ schema $\bullet$ instance $\bullet$ mapping $\bullet$ logical data
independence $\bullet$ physical data independence
\end{keyterms}

\section{One Database, Three Descriptions}

In \chref{whynotfiles} the fees program and the results program each had their
own file and their own idea of what a student was. Putting them behind one
DBMS fixes the duplication --- but it creates a new problem immediately.

The fees clerk wants a student to be a roll number, a name and an amount
owing. The examinations officer wants a roll number, a name and a list of
marks. The person who has to make the database fast wants to know that the
student table is sorted by roll number and has an index on surname. All three
are talking about the same data and none of them should have to care about the
other two.

The answer, standardised in 1975 and unchanged since, is to describe the
database \emph{three times}.

\begin{definitionbox}[The three levels]
The \term{external level} is what a particular user or program sees: their own
tailored view of the part of the database they need, in the shape they need
it. There are many external schemas, one per group of users.

\smallskip
The \term{conceptual level} is the whole database, described once, logically:
every entity, every attribute, every relationship and every constraint ---
and nothing at all about how it is stored. There is exactly one conceptual
schema.

\smallskip
The \term{internal level} --- also called the \term{physical level} --- is how it is actually stored: files, pages, record
layouts, indexes, compression. There is exactly one internal schema.

\smallskip
A \term{mapping} connects adjacent levels, so that the DBMS can translate a
request expressed against one into an operation on the next.
\end{definitionbox}

\dsfig{%
\begin{tikzpicture}[
  lvl/.style={rounded corners=3pt, draw=#1, line width=1.1pt, fill=#1!8,
              text=#1!60!black, font=\bfseries\small, align=center,
              text width=3.3cm, minimum height=1.1cm},
  ext/.style={rounded corners=2pt, draw=secondary, line width=0.9pt,
              fill=secondary!8, text=secondary!60!black,
              font=\scriptsize\bfseries, align=center,
              text width=2.5cm, minimum height=0.85cm},
  note/.style={font=\scriptsize, text=neutral, align=left, text width=5.4cm,
               anchor=west},
  arr/.style={-{Stealth[length=2.2mm]}, primary!55, line width=0.9pt},
  map/.style={font=\scriptsize\itshape, text=accent, align=center}
]
  \node[ext] (e1) at (-2.9,3.6) {Fees clerk};
  \node[ext] (e2) at (0.0,3.6)  {Examinations};
  \node[ext] (e3) at (2.9,3.6)  {Student portal};
  \node[lvl=primary]  (c) at (0,1.7) {Conceptual schema\\{\scriptsize
       \mdseries the whole database, once}};
  \node[lvl=purplehl] (i) at (0,-0.2) {Internal schema\\{\scriptsize
       \mdseries pages, records, indexes}};
  \node[draw=neutral, line width=0.9pt, fill=neutral!8, text=neutral,
        font=\scriptsize, align=center, text width=3.3cm,
        minimum height=0.6cm] (d) at (0,-1.7) {the stored database};

  \foreach \e in {e1,e2,e3}{\draw[arr] (\e) -- (c);}
  \draw[arr] (c) -- (i);
  \draw[arr] (i) -- (d);

  \node[map] at (3.6,2.65) {external/conceptual\\mapping};
  \node[map] at (3.6,0.75) {conceptual/internal\\mapping};

  \node[note, text=secondary!60!black] at (4.9,3.6)
       {Many of these. Each sees only what it needs, shaped how it needs it.};
  \node[note, text=primary!60!black] at (4.9,1.7)
       {\textbf{Exactly one.} Every entity, attribute, relationship and
        constraint --- and nothing about storage.};
  \node[note, text=purplehl!60!black] at (4.9,-0.2)
       {Exactly one. How it is really kept. Change this and nothing above
        should notice.};
\end{tikzpicture}}%
{The ANSI/SPARC three-level architecture. The two mappings are the whole
value of the arrangement: because they exist, a change at one level can be
absorbed instead of propagated.}{threelevels}

\section{Schema and Instance}

Two words that students use interchangeably for about a fortnight, and then
cannot afford to.

\begin{definitionbox}[Schema, Instance]
A \term{schema} is the \emph{description}: the names of the tables, their
columns, the types, the constraints. It is set when the database is designed
and changes rarely.

\smallskip
An \term{instance} (or \emph{state}, or \emph{extension}) is the data
\emph{in} the database at one moment. It changes every time anybody inserts a
row.

\smallskip
The analogy worth keeping: a schema is a variable's declaration, an instance
is its value. \texttt{int count;} is the schema; \texttt{7} is the instance.
\end{definitionbox}

\begin{examplebox}[The same distinction, three ways]
\begin{tabular}{@{}L{4.0cm}L{5.0cm}L{5.0cm}@{}}
& \textbf{Schema} & \textbf{Instance} \\
Programming & \texttt{int count;} & \texttt{7} \\
A database & \relation{Student}{\relkey{roll\_no}, name, programme} &
the 4{,}000 rows in it today \\
A form & the blank form & one filled-in copy \\
\end{tabular}

\smallskip
When somebody says ``the database changed'', the first question is always
which of the two they mean, because the answers are days apart in effort.
\end{examplebox}

\section{Data Independence, Which Is the Point of All This}

Three descriptions and two mappings are a lot of machinery. It buys exactly
one thing, and the thing is worth it.

\begin{definitionbox}[Logical and physical data independence]
\term{Physical data independence} is the ability to change the \emph{internal}
schema --- how data is stored --- without changing the conceptual schema or
any application program. Adding an index, moving to a different disk,
reorganising the file: none of these should require a line of application code
to change.

\smallskip
\term{Logical data independence} is the ability to change the
\emph{conceptual} schema --- adding a table, adding a column, splitting one
table into two --- without changing the external schemas or the programs that
use them.

\smallskip
Physical independence is easy and essentially achieved. \textbf{Logical
independence is hard and only partly achieved}, which is why it is worth
understanding where it fails.
\end{definitionbox}

\begin{conceptbox}[Which level does this change belong to?]
A question worth being able to answer instantly, because the cost of a change
is decided by its level.

\begin{tight}
  \item \textbf{``Add an index on surname.''} Internal. Nothing else changes.
        \chref{indexes} does exactly this and measures it.
  \item \textbf{``Add a telephone number to students.''} Conceptual. Existing
        views that never mentioned it carry on working --- this is logical
        independence succeeding.
  \item \textbf{``Split the student table into student and student-contact.''}
        Conceptual, and this is where logical independence is \emph{tested}:
        \chref{views} shows how a view can hide the split from programs that
        still want one table. \chref{normalforms} does this split for real.
  \item \textbf{``The fees clerk should not see home addresses.''} External. A
        view, and in a server DBMS a privilege (\chref{project}).
  \item \textbf{``Store dates as text rather than numbers.''} Looks internal.
        It is \emph{conceptual}, because it changes what comparisons mean ---
        and \chref{ddl} measures what SQLite does about that.
\end{tight}

\smallskip
The last one is the one that catches people, and it catches them because
``how it is stored'' sounds internal in English and is not.
\end{conceptbox}

\section{Where SQLite's Three Levels Are}

An architecture standardised in 1975 is worth nothing to you unless you can
point at its three levels in the tool on your own machine, which at first
glance appears to have only one.

\begin{alertbox}[All three are there; only two can be enforced]
It is worth knowing how an architecture from 1975 maps onto the tool in front
of you, because at first glance SQLite seems to have only one level.

\begin{tight}
  \item \textbf{Conceptual:} your \sqlkw{CREATE TABLE} statements. Exactly one
        set of them, and \texttt{.schema} prints it.
  \item \textbf{Internal:} the page structure, the B-trees, the indexes.
        \chref{firstdatabase} measures it directly and you never write it ---
        SQLite decides, and \sqlkw{PRAGMA} is the only window onto it.
  \item \textbf{External:} views (\chref{views}). SQLite has them, so the
        level exists --- but it has no users, so it cannot restrict anybody
        \emph{to} a view. The level is present and the enforcement is not,
        which is \chref{project}'s subject.
\end{tight}

\smallskip
\textbf{This is a real limitation and worth stating now rather than in week
sixteen.} You will build all three levels in this course. You will only be
able to \emph{enforce} two of them.
\end{alertbox}

\begin{pitfall}
\begin{tight}
  \item \textbf{Confusing schema with instance.} A schema change is a project;
        an instance change is a Tuesday.
  \item \textbf{Thinking there are many conceptual schemas.} There is one. Many
        external, one conceptual, one internal.
  \item \textbf{Believing physical independence means performance does not
        matter.} It means you can change storage \emph{without rewriting
        programs}. It does not mean the choice is free --- \chref{indexes}
        measures exactly how unfree it is.
  \item \textbf{Calling a storage-format change internal when it changes
        meaning.} If a comparison starts giving different answers, it was
        conceptual.
  \item \textbf{Assuming logical independence is automatic.} It is the hard
        one, and \chref{views} is how it is bought.
\end{tight}
\end{pitfall}

\begin{lab}[Find the three levels in a database you did not design]
\begin{lstlisting}[language=bash]
# --- 1. Make a database with something in it -------------------------
sqlite3 arch.db
\end{lstlisting}
then, in the shell:
\begin{lstlisting}[language=sql]
-- --- 2. State a conceptual schema -------------------------------
CREATE TABLE student (
  roll_no   TEXT PRIMARY KEY,
  name      TEXT NOT NULL,
  programme TEXT NOT NULL,
  semester  INTEGER NOT NULL
);
INSERT INTO student VALUES
  ('2021-IT-001','Ayesha Khan','BSIT',2),
  ('2021-IT-002','Bilal Ahmed','BSIT',2),
  ('2021-CS-007','Hira Nadeem','BSCS',4);

-- --- 3. Read the conceptual schema back -------------------------
.schema student
-- That text IS the conceptual schema. Nothing in it mentions storage.

-- --- 4. Build an external schema --------------------------------
CREATE VIEW fees_view AS
  SELECT roll_no, name FROM student;
SELECT * FROM fees_view;
-- A user given only this cannot see "programme" or "semester".

-- --- 5. Look at the internal schema -----------------------------
PRAGMA page_size;
PRAGMA page_count;
SELECT name, type FROM sqlite_master;
-- You did not write any of this and cannot usefully change it.

-- --- 6. Change the internal level and prove nothing else moved --
CREATE INDEX idx_name ON student(name);
SELECT * FROM fees_view;
-- Same answer. That is physical data independence, demonstrated.

-- --- 7. Now change the conceptual level -------------------------
ALTER TABLE student ADD COLUMN email TEXT;
SELECT * FROM fees_view;
-- Still works, and still has two columns. Logical independence.

-- --- 8. Break it ------------------------------------------------
-- Try to drop the column the view depends on:
ALTER TABLE student DROP COLUMN name;
-- Write down exactly what happens, and which level failed.
\end{lstlisting}

\smallskip
\textbf{What to notice.} Steps 6 and 7 are the two kinds of independence, each
demonstrated in one line. Step 8 is the limit of the second one, and the error
message is worth reading twice.
\end{lab}

\begin{checkpoint}
\begin{tightnum}
  \item Classify each of these by level, and say which kind of data
        independence (if any) protects the programs from it: adding an index;
        adding a column; renaming a column; storing the database on a faster
        disk; giving the library staff a view with no home addresses.
  \item Explain the difference between a schema and an instance to somebody who
        has written programs but never used a database, using an example that
        is not a database.
  \item SQLite has views but no users. Say precisely which of the three levels
        is affected by that, what you can still do, what you cannot, and how
        you would achieve the same effect in an application built on SQLite.
\end{tightnum}
\end{checkpoint}

\begin{chaptersummary}
\begin{tight}
  \item The ANSI/SPARC architecture describes one database three times:
        \emph{many} external schemas, \emph{one} conceptual schema, \emph{one}
        internal schema, joined by two mappings.
  \item The external level is what a user sees; the conceptual level is the
        whole database logically and says nothing about storage; the internal
        level is how it is really kept.
  \item A \emph{schema} is the description and changes rarely; an
        \emph{instance} is the data and changes constantly. A schema is a
        declaration, an instance is a value.
  \item \textbf{Physical data independence} lets the internal schema change
        without touching programs. It is essentially achieved.
  \item \textbf{Logical data independence} lets the conceptual schema change
        without touching external schemas. It is hard, partial, and bought
        with views (\chref{views}).
  \item The level a change belongs to decides what it costs. ``Store dates as
        text'' sounds internal and is conceptual, because it changes what
        comparisons mean.
  \item In SQLite the conceptual level is your \sqlkw{CREATE TABLE}
        statements, the internal level is pages and B-trees you never write,
        and the external level is views --- which exist but cannot be
        enforced, because there are no users to restrict.
\end{tight}
\end{chaptersummary}

\begin{reviewq}
  \item \textbf{(MCQ)} How many conceptual schemas does a database have?
        (a)~one per user; (b)~one per table; (c)~exactly one; (d)~one per
        application.
  \item \textbf{(MCQ)} Adding an index without changing any program
        demonstrates: (a)~logical data independence; (b)~physical data
        independence; (c)~an external schema; (d)~an instance change.
  \item \textbf{(Short)} Give one change that logical data independence
        protects you from and one that it does not, and say what makes the
        difference.
  \item \textbf{(Short)} Why does the architecture need \emph{mappings} as well
        as schemas? What would be lost without them?
  \item \textbf{(Short)} A colleague says the internal level does not matter
        because of physical data independence. Reply in two sentences.
  \item \textbf{(Applied)} A university wants three groups to use one student
        database: admissions, examinations, and the students themselves. Design
        the external schemas --- what each sees and what each must not --- and
        say which parts of your design SQLite can enforce and which it cannot.
  \item \textbf{(Applied)} Your conceptual schema must change: the single
        \texttt{address} column becomes four columns. Describe how you would
        make that change without breaking existing programs, what the view
        would look like, what it would cost, and at what point you would decide
        the view is no longer worth keeping.
\end{reviewq}
